\section*{IN2. Diseña una propuesta de investigación}

\subsection*{Texto literal del indicador}
\begin{quote}
Diseña una propuesta de investigación.
\end{quote}

\subsection*{Respuesta del grupo}
Se diseñó una revisión de literatura tipo \textit{summary paper}, documentada en la Sección II del informe.

\begin{itemize}[leftmargin=*]
  \item \textbf{Preguntas guía:} (1) ¿qué riesgos introduce la computación cuántica?, (2) ¿qué soluciones y estándares existen?, (3) ¿qué dificultades técnicas presenta la migración? y (4) ¿qué evidencia existe de adopción actual?
  \item \textbf{Fuentes de búsqueda:} Consensus para literatura académica y búsquedas web para documentos técnicos y casos recientes. Gemini se usó solo para sugerir posibles fuentes; cada sugerencia se verificó en la fuente original.
  \item \textbf{Cadenas de búsqueda:} ``post-quantum cryptography'' combinada con \textit{review}, \textit{standardization}, \textit{migration}, \textit{implementation}, \textit{performance}, \textit{NIST} y \textit{security}.
  \item \textbf{Criterios de inclusión:} autoría identificable, relación directa con PQC, evidencia verificable y aporte complementario al resto de fuentes. Además, el conjunto debía combinar las tres categorías (académica, técnica y de actualidad) con un mínimo de dos académicas.
  \item \textbf{Criterios de exclusión:} fuentes sin autoría clara, fuentes secundarias cuando existía la primaria y trabajos poco relacionados con las preguntas guía.
  \item \textbf{Alcance:} fuentes recientes en inglés y español, con énfasis en la transición organizacional y técnica, no en el diseño matemático de los algoritmos.
  \item \textbf{Componente experimental:} una demo reproducible con la biblioteca liboqs (Open Quantum Safe) para observar ML-DSA, uno de los estándares revisados.
\end{itemize}
